% chapters/chapter2/2_1_intro.tex
% !TeX root=../../main.tex

% Introduction
\section{مقدمه}

در این فصل، فرایند مدل‌سازی دینامیکی و طراحی کنترل‌کننده پایه برای سیستم جرثقیل سقفی با طول کابل متغیر تشریح می‌شود. نخست، معادلات حرکت سیستم غیرخطی با در نظر گرفتن درجات آزادی کالسکه، زاویه نوسان آونگ و مکانیزم بالابر (تغییر طول کابل) با استفاده از رهیافت تحلیلی اویلر-لاگرانژ استخراج می‌شود. پس از آن، معادلات مربوط به نیروهای تکیه‌گاهی ریل برای بررسی محدودیت‌های فیزیکی سیستم محاسبه می‌شوند.

به منظور تحلیل پایداری و طراحی کنترل‌کننده کلاسیک، مدل فضای حالت سیستم استخراج شده و حول نقطه تعادل آویزان خطی‌سازی می‌شود. در ادامه، یک کنترل‌کننده فیدبک حالت بهینه (\lr{LQR}) به عنوان معیار ارزیابی (\lr{Benchmark}) طراحی شده و رفتار سیستم تحت هدایت این کنترل‌کننده شبیه‌سازی می‌شود. تحلیل دینامیکی نتایج در این فصل، محدودیت‌های بنیادین رویکرد خطی را در مواجهه با خطاهای موقعیت بزرگ آشکار ساخته و ضرورت گذار به استفاده از برنامه‌ریز مسیر و استراتژی‌های کنترل غیرخطی در فصل‌های آتی را توجیه می‌کند.